\documentclass[../main]{subfiles}
\begin{document}
\ifSubfilesClassLoaded{\setcounter{page}{36}}{}
\section{Necessary Labour and the Conditions for Communist Activity}
\label{sec:06_necessary_labour}

In communist society, human life is not mediated by private property and the exchange of commodities; consequently, it is free from everything required to secure the conditions for the existence of private property. This gives society considerable advantages and saves its labour power. In the “Speeches in Elberfeld”, Engels offered a number of examples that show most vividly how much human labour can be saved in a society where the interests of its members are not opposed to one another and the products of activity do not turn into commodities confronting one another on the market.

The labour needed to keep the market running, the market through which social resources are distributed, is replaced by more efficient direct distribution. This makes it possible to stop performing actions that are senseless from the standpoint of organising production on a social scale, and to stop squandering people’s lifetimes on them. Producing packaging that is unnecessary for simply storing goods, carrying out shipments tied to sales rather than to delivery to the end consumer, producing goods that will later turn out to be surplus on the market and go unsold, and, still more, producing things that harm people yet bring enormous profits in a society of commodity production (drugs and the like)—none of this makes any sense in a society where production does not have a commodity character.

Abolishing the private character of appropriation also yields colossal savings of living human labour. This is the same as in the case of commodity production, but viewed from the side of the labour spent bringing products to final consumption. The examples of everyday life and transport, which we have already written about, positively cry out that it is possible to free up a vast amount of human labour that is currently expended in an inefficient, individual way.

Modern capitalism demands human effort aimed at sustaining the capacity to produce surplus value. There is more and more labour of this kind. The production of surplus value has, of necessity, been supplemented by the production of consumption and of a way of life geared towards consuming particular commodities. This alone shows the growing (and, compared with the times of Marx and Engels, greatly increased) possibilities for liberating people, who will be able, on the basis of collective property, to engage in producing what society needs for its development.

Social time is spent on maintaining the restrictions that private property imposes on the life of society and on eliminating its consequences that are harmful to society, including crime: the police, the standing army, officials, and all manner of guards of private property. All this labour power, too, can be freed up once private property in the means of production is overcome, and redirected to performing labour in which society has a direct interest.

All of this in itself, even on the basis of already existing technologies, will reduce socially necessary labour time and create room for free time to exist for genuine human life—that is, for communism proper.

Another important tendency also points to movement in the direction of communism. It is maturing within capitalism and is held back by it only because capital is the dominant social relation. “The great historic side of capital is the creation of this surplus labour, superfluous from the standpoint of mere use value, of mere subsistence of the worker, and the historic destiny of capital will be fulfilled when, on the one hand, needs have been developed to such an extent that surplus labour itself, labour beyond what is absolutely necessary for life, has become a general need arising out of people’s individual needs themselves; and when, on the other hand, general industriousness, thanks to the strict discipline of capital that successive generations have passed through, has developed as the general property of the new generation; and when, finally, this general industriousness, thanks to the development of the productive forces of labour, constantly spurred on by capital with its boundless passion for enrichment and acting under the only conditions in which it can realise this passion, has brought it about that, on the one hand, the possession and preservation of general wealth require only a comparatively small amount of labour time from society as a whole, and, on the other hand, the working society relates scientifically to the process of its progressive reproduction, its reproduction in ever-growing abundance; hence, when the kind of labour comes to an end in which a human being himself does what he can make things do for him, for man,” wrote Marx. We would add that this would mark the completion of the entire period of the existence and development of the social division of labour.

The idea of the automated factory, which emerged as early as Marx’s time, and the idea of making wide use of nature’s potential, through knowledge of it, for the good of all people, capture the really existing tendencies towards creating the conditions for communism. Communism will advance as things increasingly do the work they can do for human beings and without them, as there is less and less of the unfree human labour necessary for reproducing human life, and as the government of persons is increasingly transformed into the administration of things. Every person will be able to devote more and more time to the free play of his bodily and mental powers or, in other words, to doing what is “after his own heart”, to the benefit of society as a whole. This is precisely what constitutes the essence of communist activity.

The historical role of capital as a social relation essential to the development of the productive forces cannot be fulfilled under capitalism. Under the undivided rule of capital, capital constantly reproduces the thing-like character of labour, time and again throwing the development of the productive forces back wherever a tendency to overcome that character has emerged. For this reason, under capitalism science cannot fully become a direct force of production, and as the thing-like character of labour is reproduced, so too, among others, are its most backward forms. The complete abolition of capital as a social relation is the most important question of the period of transition to communist society. In this period, the movement of capital towards handing purely material production over to machines must be consciously carried through to the point where human activity is taken beyond the bounds of the thing-like character of labour.

Here, however, we must not forget that capital at the same time also reproduces the opposite tendency, one that leads society back to capitalism. It is this tendency, inseparably bound up with the first though it is, that socialist society can and must openly combat if the transition to communism is to take place at all. The front line between the old capitalist and the new communist objectively runs between those spheres where all the material preconditions exist for handing thing-like labour over to things once and for all, and those where this cannot yet be done without first carrying the division of labour to its utmost extreme and thereby creating the physical preconditions for transferring it to machines. Wherever the complete automation of the purely material production of things is possible, capital as a social relation, together with the corresponding character of labour, loses its relevance and therefore objectively can and must be abolished. And until this is done, capital will continue to reproduce itself as an essential social relation.

The notorious question of the anti-communist philistine, “Who will clean the toilets if everyone gets to do what they love?”, loses all meaning. Its rhetorical force is deflated not only because there is a vast difference between occasionally setting aside one’s favourite pursuit to perform this unpleasant but necessary social duty and being chained to it for life, but also because self-cleaning toilets already exist. It is only the dominance of private property, and the fact that under these conditions the labour power of cleaners, that is, the lifetime of human beings, turns out to be cheaper than things, that prevents their widespread adoption.

Under communism, by contrast, the human being, that is, the time and fullness of a person’s life, becomes the highest value, and all things matter only insofar as they secure a genuinely human life for every member of society. The very fact that a living person must spend their lifetime on uncreative labour that could be handed over to machines is intolerable for a society in which the human being is the highest value. This state of affairs must be abolished, and only then will it be possible to speak of communism developed on its own foundation.

It must be acknowledged, however, that during the period of transition to communism, necessary labour that lacks full freedom will persist for some time, until the appropriate technologies for its transcendence have been developed and introduced everywhere. But eliminating the rule of capital will make it possible to limit the working day to necessary labour, excluding labour that has meaning only within the movement of value towards its self-expansion, the results of which are time and again destroyed during crises or simply written off as “costs”. At first, necessary labour (in the sense of labour that is not yet fully free) will expand its bounds. This is needed to provide social funds and the conditions for further development. But the general direction of society’s movement towards communism, which includes the corresponding development of the productive forces, presupposes reducing such labour to ever smaller proportions. Beyond a certain point, such labour can be left out of account altogether, since it no longer has any serious influence on the general course of human life. The most important thing about necessary labour in the transition to communism, as discussed here, is that this not yet entirely free labour is labour that produces free time and a human being capable of acting freely on a social scale.
\end{document}
